% !TeX root = ../talk.tex

% Which data? One message, in the scores themselves (MusiXTeX, figures/make_musix.py): in the chorale the notes of the four
% voices fall together, in the mass each voice has its own line. A score, the name, the size of the corpus, one line on
% the writing and one on the rule. The rest of the dataset (link, source, licence, who else uses it) is in the footnotes,
% which carry only the numbers of the references.
\begin{frame}{Which data?}
    \begin{columns}[T,onlytextwidth]
        \begin{column}{0.47\textwidth}
            \centering
            \parbox[t][3.0cm][t]{\linewidth}{\centering\scalebox{0.68}{{\smallmusicsize\elemskip=2.67pt
\begin{music}
\parindent12mm\nobarnumbers
\instrumentnumber{4}
\setstaffs11\setclef1{6}\setname1{\scriptsize Bass}
\setstaffs21\setclef2{0}\setname2{\scriptsize Tenor}
\setclefsymbol2{\treblelowoct}
\setstaffs31\setclef3{0}\setname3{\scriptsize Alto}
\setstaffs41\setclef4{0}\setname4{\scriptsize Soprano}
\generalsignature{0}\generalmeter{\meterfrac44}
\startextract
\Notes \zql{5}&\zql{4}&\zcu{0}&\zqu{2}\sk\en
\Notes &&\zcu{-1}&\sk\en
\Notes \zql{6}&\zcu{3}&\zcu{-2}&\zqu{3}\sk\en
\Notes &\zcl{4}&\zcu{-1}&\sk\en
\Notes \zcu{3}&\zql{5}&\zcu{0}&\zqu{2}\sk\en
\Notes \zcl{4}&&\zcu{1}&\sk\en
\Notes \zcl{5}&\zcl{5}&\zqu{2}&\zql{5}\sk\en
\Notes \zcl{6}&\zcl{6}&&\sk\en\bar
\Notes \zql{7}&\zcl{7}&\zqu{2}&\zql{5}\sk\en
\notes &\zccl{6}&&\sk\en
\notes &\zccl{5}&&\sk\en
\NOtes \zqu{0}&\zql{6}&\zqu{2}&\zql{4}\sk\en
\NOtes \zqu{3}&\zql{7}&\zqu{2}&\zql{5}\sk\en
\NOtes \zqu{3}&\zql{7}&\zqu{2}&\zql{5}\sk\en
\zendextract
\end{music}}
}}\par
            {\Large\bfseries Bach chorales\footnotemark[1]}\par
            {\footnotesize\color{nord3}371 chorales (351 four-voice)\par}
            \medskip
            {\large Voices move together\par}
            {\color{nord3}No parallel fifths or octaves \cite{zhu2021}}
        \end{column}
        \begin{column}{0.06\textwidth}
            \centering\textcolor{nord4}{\rule{0.6pt}{5.2cm}}
        \end{column}
        \begin{column}{0.47\textwidth}
            \centering
            \parbox[t][3.0cm][t]{\linewidth}{\centering\scalebox{0.68}{{\smallmusicsize\elemskip=2.74pt
\begin{music}
\parindent12mm\nobarnumbers
\instrumentnumber{4}
\setstaffs11\setclef1{6}\setname1{\scriptsize Bass}
\setstaffs21\setclef2{0}\setname2{\scriptsize Tenor}
\setclefsymbol2{\treblelowoct}
\setstaffs31\setclef3{0}\setname3{\scriptsize Alto}
\setstaffs41\setclef4{0}\setname4{\scriptsize Soprano}
\generalsignature{-1}\generalmeter{\meterfrac42}
\startextract
\NOTEs \rlap{\pause}&\zw{6}&\rlap{\pause}&\rlap{\pause}\sk\en
\NOTes \zw{11}&\zhlp{8}&\rlap{\hpause}&\zw{3}\sk\en
\NOtes &&\zhu{-1}&\sk\en
\NOtes &\zql{9}&&\sk\en\bar
\NOTes \zw{11}&\zhl{10}&\zhup{1}&\zw{3}\sk\en
\NOtes &\zhl{8}&&\sk\en
\NOtes &&\zqu{2}&\sk\en
\NOTes \zw{8}&\zhl{7}&\zhu{3}&\zw{5}\sk\en
\NOTes &\zhl{10}&\zhu{1}&\sk\en
\zendextract
\end{music}}
}}\par
            {\Large\bfseries Palestrina\footnotemark[2]}\par
            {\footnotesize\color{nord3}1,318 movements (487 four-voice)\par}
            \medskip
            {\large Each voice has its own line\par}
            {\color{nord3}Dissonances on weak beats \cite{anders2019}}
        \end{column}
    \end{columns}
    \footnotetext[1]{\scriptsize\href{https://github.com/cuthbertLab/music21/tree/master/music21/corpus/bach}{music21/corpus/bach}. Riemenschneider edition. Public domain music, music21 terms. Also in \cite{hadjeres2017,zhu2021}.}
    \footnotetext[2]{\scriptsize\href{https://github.com/cuthbertLab/music21/tree/master/music21/corpus/palestrina}{music21/corpus/palestrina}. Le Opere Complete, enc. J. Miller, \copyright\ 2000; no licence text. Also in \cite{anders2019,padilla2018}.}
\end{frame}
